\cotefile{157-2}
\scanpage{147}
\mspage{92}
ce qui \unsure{résultera} aussitôt du corollaire :

\uline{Corollaire} \unsure{Soient|Soit} un diagramme dans $\mathcal M_{/S}$
\diagram{\begin{tikzcd}[ampersand replacement=\&]
X\arrow[r,"i"] \& C\arrow[r,shift left=.6ex,"u_1"]\arrow[r,shift right=.6ex,"u_2"'] \& C'\text{,}
\end{tikzcd}}
avec $C$ et $C'$ dans $\underline{\mathrm{Fib}}\,S$ (i.e. $C\longrightarrow S$ et $C'\longrightarrow S$ dans $\mathrm{Fib}$), et $i\in\mathrm{Cof}_0$\inserted{ et $\unsure{i|i'}\in W$}. Si $u_1i=u_2i$, alors $u_1$ et $u_2$ définissent la $\eqsim$ flèche de $(\underline{\mathrm{Fib}}\,S)W_S^{-1}$.

\uline{Dém} Soit
\[
u=(u_1,u_2):C\longrightarrow C'\times C'\text{,}
\]
qui rend commutatif par hyp. le diagramme
\diagram{\begin{tikzcd}[ampersand replacement=\&]
X\arrow[r,"w"]\arrow[d,"i"'] \& C'\arrow[d,"\mathrm{diag}_{C'}"]\\
C\arrow[r,"u"'] \& C'\times_S C'\text{,}
\end{tikzcd}}
où $w=u_1i=u_2i$. Factorisons $\mathrm{diag}_{C'}$ en
\diagram{\begin{tikzcd}[ampersand replacement=\&]
C'\arrow[r,"i'"] \& \overline{C'}\arrow[r,"\delta'"] \& C'\times_S C'\text{,}\quad\text{avec }\delta'\in\mathrm{Fib}\text{,}
\end{tikzcd}}
\deleted{\unsure{pas|par}} et $i'\in W$. On trouve un carré comm.
\diagram{\begin{tikzcd}[ampersand replacement=\&]
X\arrow[r,"w'=i'w"]\arrow[d,"i"'] \& \overline{C'}\arrow[d,"\delta'"] \& {}\\
C\arrow[ur,dashed,"\overline u"]\arrow[r,"u"] \& C'\times_S C'\arrow[r,shift left=.6ex,"\mathrm{pr}_1"]\arrow[r,shift right=.6ex,"\mathrm{pr}_2"'] \& C'
\end{tikzcd}}
avec $i\in\mathrm{Cof}_0$, $\delta'\in\mathrm{Fib}$, d'où l'existence de la flèche en pointillés, rendant les deux triangles commutatifs. On veut prouver $\mathrm{pr}_1\circ u=\mathrm{pr}_2\circ u$ dans $(\underline{\mathrm{Fib}}\,S)W_S^{-1}$, comme $u=\delta'\overline u$ il suffit de prouver que l'on a $\mathrm{pr}_1\delta'=\unsure{p_2|\mathrm{pr}_2}\delta'$ dans $(\underline{\mathrm{Fib}}\unsure{,}S)W_S^{-1}$.
% archive sheet number: 146

\scanpage{148}
\mspage{93}
Il faut seulement vérifier \inserted{d'abord} (pour justifier l'argument) que cette assertion a un sens, i.e. que $\overline{C'}$ est lui aussi \unsure{dans} $\underline{\mathrm{Fib}}_{/S}$. Mais les hyp. de stabilité de \inserted{$F={}$}$\mathrm{Fib}$\unsure{-} par \unsure{chgt.|$\partial$b.} de base et composition assurent que 1°) $C'\times_S \deleted{\unsure{S}}C'$ est dans $\underline{\mathrm{Fib}}\,S$ (qui est stable par produit), et par suite 2°) que $\overline{C'}$ l'est aussi, puisque $\delta'\in\mathrm{Fib}$.
Comme \inserted{$C'$ est aussi dans $\underline{\mathrm{Fib}}\,S$ et} $C'\xrightarrow{i'}\overline{C'}$ est dans $W_S$, pour vérifier $\mathrm{pr}_1\delta'=\mathrm{pr}_2\delta'$ dans le localisé de $\underline{\mathrm{Fib}}\,S$, il suffit de voir que $(\mathrm{pr}_1\delta')i'=(\mathrm{pr}_2\delta')i'$, or comme $\delta'i'=\mathrm{diag}_{C'}$ on a l'égalité même dans $\underline{\mathrm{Fib}}\,S$ (sans localiser).

On a donc défini $\gamma(X)$, $\gamma(u)$ pour objets et flèches de \deleted{$\unsure{\mathrm{Fib}|\mathrm{Fib}_S}$} $\mathcal M_{/S}$. La transitivité pour $\gamma(u)$ est claire, et aussi $\gamma(\mathrm{id})=\mathrm{id}$. Donc on tient le foncteur
\[
\gamma:\mathcal M_{/S}\longrightarrow(\underline{\mathrm{Fib}}\,S)W_S^{-1}
\]

Soit maintenant $\unsure{\mathcal F|F}$ une sous-catégorie pleine de $\mathcal M_{/S}$, avec
\[
\underline{\mathrm{Fib}}\,S\subset\mathcal F\subset\mathcal M_{/S}\text{,}
\]
considérons $\gamma$ comme foncteur à val. dans $\mathcal F W_S^{-1}$, (j'ai oublié de vérifier que $\gamma$ transforme $W_S$ \unsure{de|ds} $W_S$, mais c'est tautologique.) Ainsi on a des foncteurs
\diagram{\begin{tikzcd}[ampersand replacement=\&]
\mathcal F W_S^{-1}\arrow[r,shift left=.6ex,"\delta"] \& (\mathcal M_{/S})W_{\unsure{S|\mathcal F}}^{-1}\arrow[l,shift left=.6ex,"\gamma"]\text{,}
\end{tikzcd}}
il faut voir qu'ils sont quasi-inv. l'un de
% archive sheet number: 147

\scanpage{149}
\mspage{94}
l'autre.

Du côté de $\mathcal M_{/S}$ : on doit comparer $X$ et $\delta\gamma(X)=C(X)$, c'est fait par
\[
i_X:X\longrightarrow C(X)
\]
qui est dans $\unsure{W_S|W}$, d'où l'isom. dans le localisé :
\[
X\simeq\delta\gamma(X)\quad\text{dans }(\mathcal M_{/S})W_S^{-1}\text{.}
\]
Il faut vérifier que c'est fonctoriel en $X$.

Du côté de $\unsure{\mathcal F|F}$ \deleted{$\underline{\mathrm{Fib}}\,S$}, comparer $X$ et $\gamma\delta(X)$, i.e. $X$ et $C(X)$ : Même tabac, on choisit $i_X:X\simeq C(X)=\gamma\delta(X)$. À vérifier la fonctorialité.
\figure{\begin{tikzpicture}[x=1cm,y=1cm]\draw (0,.16) .. controls (.02,.10) and (.09,.04) .. (.14,0);\end{tikzpicture}}

Il s'agit maintenant d'introduire un localisé plus fin de $\mathcal M_{/S}$. Pour ceci, on suppose \unsure{donné|donnée} une autre classe de morphismes, je l'appelle $\mathcal L$ (comme ``localisation'')
\[
\mathcal L\subset\mathrm{Fl}(\mathcal M)\text{,}
\]
p. ex. dans le cas de $\mathcal M=\mathrm{Cat}$, $\mathcal L$ \deleted{\unsure{sera}} sera formé des \deleted{\unsure{co}} isom. locaux$^{*}$ (s'il s'agit de définir $W_S^{\mathrm g}$), ou des \unsure{coisom.} locaux$^{**}$ (s'il s'agit de définir $W_S^{\mathrm d}$). On suppose, \inserted{que $\mathcal L$ contient les iso, et} $\mathcal L$ stable par composition, \deleted{\unsure{($^{***}$) qu'il contient les iso}}\,\unsure{,} mais non par chgt. de base.
% slanted
\margin{($^{*}$) ou aussi, des morphismes \uline{lisses}. (ça doit donner le $\eqsim$ résultat)}
% slanted
\margin{($^{**}$) ou aussi, des morphismes de projection}
% slanted
\margin{$^{***}$ Il n'y en \unsure{est|ait} peut-être pas besoin.}
On définit alors
\[
W_S^{\mathcal L}\subset\mathrm{Fl}(\mathcal M_{/S})
\]
par
\[
f\in W_S^{\mathcal L}\overset{\mathrm{def}}{\Longleftrightarrow}(f\times_S S')\in W,\quad\forall\,(S'\longrightarrow S)\in\mathcal L\text{.}
\]
% archive sheet number: 148

\scanpage{150}
\mspage{95}
\deleted{Il faut supposer que $\mathcal L$ est assez \unsure{abondant} pour}
Comme \deleted{$\unsure{F|S}$} $\mathcal L$ contient $\mathrm{id}_S$, on aura
\[
W_S^{\mathcal L}\subset W_S\text{,}
\]
d'où un foncteur canonique
\diagram{\begin{tikzcd}[ampersand replacement=\&]
(\mathcal M_{/S})(W_S^{\mathcal L})^{-1}\arrow[r,"\rho"]\arrow[d,equals,"\mathrm{def}"'] \& \underline{(\mathcal M_{/S})W_S^{-1}}\arrow[d,equals,"\mathrm{def}"]\text{,}\\
\mathrm{Ho}^{\mathcal L}_W(S) \& \mathrm{Ho}^{\mathrm{lc}}_W(S)
\end{tikzcd}}
qui est un foncteur de localisation. Je voudrais définir un foncteur en sens inverse : Pour ceci, il \deleted{faut} \inserted{est naturel de} supposer (en analogie avec la situation de $\mathrm{Cat}$)

\margin{\tikz{\draw (0,0)--(0,-1.5); \draw (.07,0)--(.07,-1.5);}}\uline{Hyp 1} Sur $\underline{\mathrm{Fib}}\,S$, $W_S=W_S^{\mathcal L}$, i.e. si $f:X\longrightarrow Y$ est dans $W$ et $X,Y$ dans $\underline{\mathrm{Fib}}\,S$, alors \deleted{$\unsure{f}\ \forall(S'\longrightarrow S)$,} pour tt chgt. de base $S'\longrightarrow S$ qui est dans $\mathcal L$, $f':X_{S'}\longrightarrow Y_{S'}$ est encore dans $W$.

Dans le contexte de Quillen, cette hypothèse \deleted{\unsure{résultera} des axiomes de Qu. quel que} a une nette tendance à être vérifiée quel \deleted{soit} que soit $\mathcal L$, i.e. pour \uline{tt} chgt. de base $S'\longrightarrow S$ sans restriction. Il suffit en fait que la structure de Quillen soit ``propre'', plus précisément (et moins restrictivement), \unsure{qu'un|qu'une} \deleted{tout changement de} $\in W$ reste tel après chgt. de base par un morphisme $\in\mathrm{Fib}$.
Dégageons ceci comme un lemme

\uline{Lemme} L'hypothèse est vérifiée \deleted{\unsure{si on|sans}} \unsure{suppose|supposé|suppos}
% archive sheet number: 149

\scanpage{151}
\mspage{96}
que $F=\mathrm{Fib}$ et $W$ sont reliés par la condition supplémentaire (en plus de la factorisation $i_X$) : pour tt carré cartésien
\diagram{\begin{tikzcd}[ampersand replacement=\&]
X\arrow[d,dash,"f"'] \& X'\arrow[l,"q" description]\arrow[d,"f'"]\\
S \& S'\arrow[l,"p" description]\text{,}
\end{tikzcd}}
Si $f$ est dans $\mathrm{Fib}$ et $p$ dans $W$, alors $q\in W$.

\uline{Dém}. Soit
\diagram{\begin{tikzcd}[ampersand replacement=\&]
X\arrow[rr,"f"]\arrow[dr,dash,"p"'] \& {} \& Y\arrow[dl,dash,"q"]\\
{} \& S \& {}
\end{tikzcd}}
avec $p,q\in\mathrm{Fib}$, $f\in W$, montrons que pour \uline{tout} chgt de base $g:S'\longrightarrow S$, $f_{S'}:X_{S'}\longrightarrow Y_{S'}$ est encore dans $W$. On factorise $g$ en
\diagram{\begin{tikzcd}[ampersand replacement=\&]
S'\arrow[r,"i"]\arrow[rr,bend right=16,"g"'] \& \overline{S'}\arrow[r,"\overline g"] \& S\text{,}\qquad i\in W,\ \unsure{g|\overline g}\in\mathrm{Fib}\text{,}
\end{tikzcd}}
et on est ramené à prouver que l'assertion est vraie dans les deux cas particuliers où $g\in\mathrm{Fib}$, ou $g\in W$. Le cas $g\in\mathrm{Fib}$ résulte aussitôt de l'hyp. du lemme (appliquée au \deleted{chgt de base} carré cartésien
\deleted{\diagram{\begin{tikzcd}[ampersand replacement=\&]
X\arrow[r,"f"]\arrow[d,"\unsure{p|q}"'] \& Y\arrow[d,"p\in\mathrm{Fib}"]\\
X'\arrow[r,"f'"'] \& Y'
\end{tikzcd}}}
\diagram{\begin{tikzcd}[ampersand replacement=\&]
X'\arrow[r,"f'"]\arrow[d,"q"'] \& Y'\arrow[d,"p"',"{\deleted{\unsure{g}\in\mathrm{Fib}}}"]\\
X\arrow[r,"f"'] \& Y\text{,}
\end{tikzcd}}
\deleted{\unsure{D}} \deleted{\unsure{le cas}} où $p\in\mathrm{Fib}$.
\inserted{Dans le cas} où $S'\longrightarrow S$ est dans $W$, il s'ensuit par l'hyp. du lemme appliquée à
\diagram{\begin{tikzcd}[ampersand replacement=\&]
Y\arrow[d,"{\deleted{\unsure{q}}}"',"\in\mathrm{Fib}"] \& Y'\arrow[l,"p" description]\arrow[d]\\
S \& S'\arrow[l]
\end{tikzcd}}
que :
% archive sheet number: 150

\scanpage{152}
\mspage{97}
$p:Y'\longrightarrow Y$ est $\in W$, et de $\eqsim$ pour $q:X'\longrightarrow X$, et $f\in W$ implique \unsure{encore|aussi|donc} $f'\in W$.
\figure{\begin{tikzpicture}[x=1cm,y=1cm]\draw[line width=.7pt] (0,.18)--(.15,0);\end{tikzpicture}}

Revenons \unsure{à|:} l'hyp. p. 95. On aura donc un foncteur
\[
(\underline{\mathrm{Fib}}\,S,W_S=W_S^{\mathcal L})\longrightarrow(\deleted{\unsure{C}}\mathcal M_{/S},W_S^{\mathcal L})
\]
d'où pour les localisés un foncteur horizontal
\diagram{\begin{tikzcd}[ampersand replacement=\&]
(\underline{\mathrm{Fib}}\,S)W_S^{-1}\arrow[r,"i_0"]\arrow[dd,"\approx"] \& (\mathcal M_{/S})(W_S^{\mathcal L})^{-1}\arrow[d,equals,"\mathrm{def}"]\arrow[ddl]\\
{} \& \mathrm{Ho}^{\mathcal L}_W(S)\\
\mathrm{Ho}^{\mathrm{lc}}_W(S)\overset{\mathrm{def}}{=}(\mathcal M_{/S})W_S^{-1} \& {}
\end{tikzcd}}
qui factorise le foncteur vertical induit par l'inclusion $(\underline{\mathrm{Fib}}\,\deleted{\unsure{S}}S,W_S)\longrightarrow(\deleted{\mathrm{Cat}}\mathcal M_{/S},W_S)$ --- foncteur dont on a vu (dans les deux cas envisagés) que c'est une équivalence.
En inversant ce dernier, \inserted{et le composant avec $i_0$,} on trouve donc un foncteur
\margin{\unsure{*}}
\diagram{\begin{tikzcd}[ampersand replacement=\&]
\mathrm{Ho}^{\mathrm{lc}}_W(S)\arrow[r,"i"]\arrow[d,equals,"\mathrm{def}"'] \& \mathrm{Ho}^{\mathcal L}_W(S)\arrow[d,equals,"\mathrm{def}"]\\
(\mathcal M_{/S})W_S^{-1} \& (\mathcal M_{/S})\deleted{\unsure{(}}(W_S^{\mathcal L})^{-1}\text{,}
\end{tikzcd}}
tel que $\rho i\simeq\mathrm{id}$, donc une \uline{rétraction} de $\mathrm{Ho}^{\mathcal L}_W(S)$ sur $\mathrm{Ho}^{\mathrm{lc}}_W(S)$. On s'attend que $i$ soit pl fidèle, et plus précisément encore,
% archive sheet number: 151

\scanpage{153}
\mspage{98}
que $\rho$ et $i$ soient adjoints, \deleted{Mais le cas de $\mathrm{Cat}$ \inserted{(si je ne m'y suis \unsure{trompé})} \unsure{nous montre} qu'ils peuvent être adjoints} $i$ étant l'adjoint à droite de $\rho$ :
\diagram{\begin{tikzcd}[ampersand replacement=\&]
\mathrm{Ho}^{\mathrm{lc}}_W(S)\arrow[r,bend right=12,"i"'] \& \mathrm{Ho}^{\mathcal L}_W(S)\arrow[l,bend right=12,"\rho"']\text{,}
\end{tikzcd}}
i.e on cherche
\[
\begin{aligned}
i\rho &\xleftarrow{\overset{\text{\uline{pas} iso}}{\unsure{i|\widetilde{\imath}}}}\mathrm{id} &&\text{dans }\mathrm{Ho}^{\mathcal L}_W &&(\text{loc. par }W_S^{\mathcal L})\\
\rho i &\xrightarrow[p]{\mathrm{iso}}\mathrm{id} &&\text{dans }\mathrm{Ho}^{\mathrm{lc}}_W &&(\text{loc par }W_S)
\end{aligned}
\]
donc pour $X$ dans $\mathcal M_{/S}$, on cherche
\[
\begin{aligned}
\unsure{i_X|\widetilde{\imath}_X}:&\quad X\xrightarrow{\deleted{\unsure{\rho i}}} C(X) &&\substack{\text{tout au moins}\\\text{dans }(\mathcal M_{/S})(W_S^{\mathcal L})^{-1}}\\
p_X:&\quad X\xleftarrow{\deleted{\unsure{p}}} C(X) &&\substack{\text{tout au moins dans}\\(\mathcal M_{/S})(W_S)^{-1}}\text{.}
\end{aligned}
\]
On choisira pour $i_X$ la flèche ``canonique'' de $X$ vers $C(X)$, définie dans $\mathcal M_{/S}$ lui-$\eqsim$, sans avoir à localiser. Mais dans le ``cas Quillen'' (quand $C$ n'est pas donné comme foncteur $\mathcal M_{/S}\longrightarrow\underline{\mathrm{Fib}}\,S$), il faut voir que cet $i_X$ est fonctoriel en $X$, quand on le regarde dans le localisé $(\mathcal M_{/S})(W_S)^{-1}$. Mais c'est contenu dans la construction de $C$ comme \uline{foncteur}
\[
\mathcal M_{/S}\longrightarrow(\underline{\mathrm{Fib}}\,S)(W_S)^{-1}
\]
On notera que $i_X$ est dans $W_S$ sans plus, pas dans $W_S^{\mathcal L}$, donc ne définit pas en général un iso de $\mathrm{Ho}^{\mathcal L}_W$.

On définit $p_X$ à l'aide de $i_X$, comme l'inverse de $i_X$ \uline{dans} $(\mathcal M_{/S})(W_S)^{-1}$, qui existe car
% archive sheet number: 152

\scanpage{154}
\mspage{99}
$i_X$ est dans $W_S$, donc iso dans $(\mathcal M_{/S})(W_S)^{-1}=\mathrm{Ho}^{\mathrm{lc}}_W$, et $p$ est donc iso. On retrouve donc $\rho i\simeq\mathrm{id}$, donc que $\rho$ est une \uline{rétraction} de $\mathrm{Ho}^{\mathcal L}_W$ sur $\mathrm{Ho}^{\mathrm{lc}}_W$.

Pour prouver l'adjonction, en procédant comme dans le cas de $\mathrm{Cat}$, on est ramené à prouver que pour tt $X$,
\[
\left\{\begin{aligned}
&C(i_X)=i_{C(X)}:\quad C(X)\rightrightarrows C\deleted{\unsure{C|(}}(C(X))\\
&\text{\uline{dans }}\mathrm{Ho}^{\mathcal L}_W\qquad(\text{loc. par }W_S^{\mathcal L})\text{.}
\end{aligned}\right.
\]

J'y étais arrivé dans $\mathrm{Cat}$, en utilisant un morphisme
\[
C(C(X))\xrightarrow{c}C(X)
\]
(essentiellement, un morphisme de composition des chemins) tel que $c\in W_S^{\mathcal L}$, et que $c\circ C(i_X)=c\circ i_{C(X)}$. Ça avait l'air très délicat, car c'était à grand peine que j'étais arrivé à construire un foncteur $\underline{\mathrm{Ch}}$ (``chemins'') \unsure{,} raisonnable, où on puisse bel et bien composer les chemins ! Je renonce pour le moment à axiomatiser les propriétés du foncteur $\underline{\mathrm{Ch}}$ dont j'ai eu besoin pour faire marcher la machine (je devrais y revenir tantôt).
Je vais commencer plutôt par le ``cas Quillen'', avec des $i_X\in\mathrm{Cof}_0$.
% archive sheet number: 153

\scanpage{155}
\mspage{100}
On reprend la description du ``foncteur'' $C$ (à valeurs dans $(\underline{\mathrm{Fib}}\,S)W_S^{-1}$, et alors sans guillemets), on trouve que $C(i_X)$, où
\[
X\xrightarrow{u=i_X}X'=C(X)
\]
est défini simplement (dans $(\underline{\mathrm{Fib}}_S)W_S^{-1}$) par ``la'' flèche $C(u)$ rendant commutatif le diagramme (dans $\underline{\mathrm{Fib}}\,S$)
\diagram{\begin{tikzcd}[ampersand replacement=\&]
X\arrow[r,"{u\,(=i_X)}"]\arrow[d,"i_X"'] \& \deleted{\unsure{C}}X'\,(=C(X))\arrow[d,"{i_{X'}\,(=i_{C(X)})}"]\\
C(X)\arrow[r,dash,"{C(u)(=C(i_X))}"] \& C(X')\,(=C(C(X))\text{,}
\end{tikzcd}}
de sorte qu'on a dans $\deleted{\unsure{F}}\mathcal M_{/S}$ le diagramme
\diagram{\begin{tikzcd}[ampersand replacement=\&]
X\arrow[r,shift left=.7ex,"{\substack{\mathrm{Cof}_0\cap W\\\unsure{\swarrow|\ni}\\i_X}}"]\arrow[r,shift right=.7ex,phantom,"{\deleted{\unsure{\longleftarrow|\longrightarrow}}}" description] \& \overset{\substack{\mathrm{Ob}\,\underline{\mathrm{Fib}}\,S\\\unsure{\swarrow|\ni}}}{C(X)}\arrow[r,shift left=.6ex,"i_{C(X)}"]\arrow[r,shift right=.6ex,"C(i_X)"'] \& \overset{\substack{\mathrm{Ob}\,\underline{\mathrm{Fib}}\,S\\\unsure{\downarrow|\ni}}}{C(C(X))}\text{,}
\end{tikzcd}}
commutatif i.e. $i_{C(X)}i_X=C(i_X)i_X$. Le Cor. page 92 nous dit que les flèches $i_{C(X)}$ et $C(i_X)$ sont égales dans $(\underline{\mathrm{Fib}}\,S)W_S^{-1}$.
Mais par hypothèse (p. 95), dans $\underline{\mathrm{Fib}}\,S$ on a $W_S=W_S^{\mathcal L}$, donc les flèches sont égales dans $(\underline{\mathrm{Fib}}\,S)(W_S^{\mathcal L})^{-1}$, et a fortiori dans $(\mathcal M_{/S})(W_S^{\mathcal L})^{-1}$, et on gagne.

Voyons ce qu'on peut faire dans le cas où on dispose d'un \uline{foncteur} $C$. On supposera bien sûr que ce foncteur dans $\mathcal M_{/S}$
% archive sheet number: 154

\scanpage{156}
\mspage{101}
n'est pas lié au $S$ particulier de $\mathcal M$ qu'on a choisi comme base, que sa construction provient à la \unsure{Cuba-Libre|Cuba-Serre|Cuba-Swan} d'un foncteur
\[
X\longmapsto\mathrm{ch}(X):\quad\mathcal M\longrightarrow\mathcal M\text{,}
\]
(``objet des chemins'') et de morphismes fonctoriels
\diagram{\begin{tikzcd}[ampersand replacement=\&]
X\arrow[r,"\alpha_X"] \& \mathrm{Ch}(X)\arrow[r,"p_X"] \& X\times X\text{,}\quad\boxed{p_X\alpha_X=\mathrm{diag}_X}
\end{tikzcd}}
% slanted
\margin{\uline{NB} On pose $\mathrm{pr}_1\circ p_X=\sigma_X$, $\mathrm{pr}_2\circ p_X=\beta_X$ (flèches ``source'' et ``but'' pour l'objet $X$).}
\unsure{précisons|prévisibles} qu'il faudra supposer \unsure{seulement} ceci.

\uline{Hyp 2}\quad $\alpha_X\in W,\quad p_X\in\mathrm{Fib}$.

Pour un $X$ sur $S$, on construira \unsure{((} $C_S(X)$ \inserted{ou $C(X)$} ainsi :
\diagram{\begin{tikzcd}[ampersand replacement=\&]
\deleted{\unsure{C(X)}} \& \deleted{\unsure{S}}\mathrm{Ch}(S)\arrow[d,"p_S"]\arrow[dl,dashed,"\beta_S" description]\arrow[dd,bend right=16,dashed,"\sigma_S"']\arrow[dr,phantom,"\mathrm{cart}" description] \& C_S(X)=C(X)\arrow[l,"\pi"']\arrow[d,"p'"]\arrow[dll,bend right=43,"{\overline f_X=f_{C(X)}}" description]\\
S \& S\times S\arrow[l,"\mathrm{pr}_2^S"']\arrow[d,"\mathrm{pr}_1^S"]\arrow[dr,phantom,"\mathrm{cart}" description] \& X\times S\arrow[l,"{\pi_1=f_X\times\mathrm{id}_S}"]\arrow[d,"\mathrm{pr}_1"]\\
\deleted{\unsure{X}} \& S\arrow[l,phantom,"{\deleted{\xleftarrow{f_X}}}" description] \& X\,(=S')\arrow[l,"f_X"']\arrow[uu,bend right=31,dashed,"i_X"']
\end{tikzcd}}
Il est l'image inverse de $(\mathrm{Ch}(S),\,\mathrm{pr}_1^S p_S=\sigma_S)$ par $X\xrightarrow{f_X}S$, mais on a explicité la factorisation de $\sigma_S$ sur le diagramme, pour réinterpréter la flèche structurale $\overline f_X=f_{C(X)}$, définie comme
\[
f_{C(X)}=\beta_S\pi\text{,}\qquad(=\mathrm{pr}_2^S\underbrace{p_S\pi}_{\substack{\shortparallel\\\pi_1p'}})
\]
par la formule
\[
f_{C(X)}=\underbrace{(\mathrm{pr}_2^S\pi_1)}_{\substack{\text{projection}\\\mathrm{pr}_2\text{ de }X\times S\\\text{dans }S}}\,p'\qquad\unsure{\checkmark}
\]
% archive sheet number: 155

\scanpage{157}
\mspage{102}
Or
\[
\begin{aligned}
&p'\in\mathrm{Fib}, &&\substack{\text{car déduit de }p_S\in\mathrm{Fib}\text{ par}\\\text{chgt de base }\pi_1,}\\
&(\mathrm{pr}_2:X\times S\longrightarrow S)\in\mathrm{Fib} &&\substack{\text{pourvu que }X\longrightarrow e\\\text{soit dans }\mathrm{Fib}\quad(\eqsim\text{ raison})}
\end{aligned}
\]
\margin{$\swarrow$ objet final de $\mathcal M$}
donc si $X\in\underline{\mathrm{Fib}}\,e$, on a $\overline f_{\deleted{\unsure{X}}C(X)}\in\mathrm{Fib}$. Comme c'est une propriété \unsure{essentielle|essentiels} qu'on veut avoir pour \uline{tout} $X$ dans $\mathcal M_{/S}$, on va faire l'hypothèse supplémentaire

\uline{Hyp 3}\quad $\forall X\in\mathrm{Ob}\,\mathcal M$, $X\longrightarrow e$ est dans $\mathrm{Fib}$.

Quant à \deleted{\unsure{la}}
\[
i_X:X\longrightarrow C(X)\text{,}
\]
il est déduit de la section $\alpha_S$ de $(\mathrm{Ch}(S),\inserted{\sigma_S)}$ sur $S$ par changement de base $X\longrightarrow S$. Notons maintenant que \unsure{G|S}
\diagram{\begin{tikzcd}[ampersand replacement=\&]
\mathrm{Ch}(S)\arrow[r,shift left=.6ex,"\sigma_S"]\arrow[r,shift right=.6ex,"\beta_S"'] \& S\qquad\text{sont dans }W\text{,}
\end{tikzcd}}
car leurs composés avec $S\xrightarrow{\alpha_S}\mathrm{Ch}(S)$ sont l'identité, et $\alpha_S\in W$. D'autre part, $\sigma_S=\mathrm{pr}_1^S\circ p$ est dans $\mathrm{Fib}$, comme composé de deux éléments de $\mathrm{Fib}$, (NB $S\longrightarrow e$ dans $\mathrm{Fib}$ par Hyp. $\unsure{3|4}$). Ainsi $\sigma_S,\beta_S$ sont dans $\mathrm{Fib}\cap W$.

On aura besoin de

\uline{Hyp 4}\quad \deleted{\unsure{L'élément|les éléments}} \inserted{La \unsure{classe}} de flèches $\mathrm{Fib}\cap W$ est stable par tt chgt. de \deleted{base}
% archive sheet number: 156

\scanpage{158}
\mspage{103}
Cela implique donc que
\[
C(X)\xrightarrow[\mathrm{pr}_1\pi']{}X
\]
est $\in W$ (étant déduit de $\sigma_S:\mathrm{Ch}(S)\longrightarrow S$ par chgt. de base), donc \deleted{toute} section \inserted{$i_X:X\longrightarrow C(X)$} de $C(X)$ sur $X$ est dans $i_X$. On a donc bien une factorisation, fonctorielle en $X$
\margin{\unsure{''}}
\diagram{\begin{tikzcd}[ampersand replacement=\&]
{} \& C(X)\arrow[dr,"{\overline f_X=f_{C(X)}}"] \& {}\\
X\arrow[ur,"i_X"]\arrow[rr,"{f_X\deleted{\unsure{X}}}"] \& {} \& S
\end{tikzcd}}
avec
\[
i_X\in W,\quad f_{C(X)}\in\mathrm{Fib}\text{.}
\]
C'est tout ce qu'on avait utilisé pour établir
\[
(\underline{\mathrm{Fib}}\,S)W_S^{-1}\mathrel{\overset{\approx}{\rule{1.7em}{.4pt}}}(\mathcal M_{/S})W_S^{-1}\text{.}
\]

Mais il s'agit maintenant, sous l'hyp. où on introduit un $\mathcal L$, de prouver que pour tt $X$ dans $\mathcal M_{/S}$, les deux flèches
\diagram{\begin{tikzcd}[ampersand replacement=\&]
C(X)\arrow[r,shift left=.6ex,"C(i_X)"]\arrow[r,shift right=.6ex,"i_{C(X)}"'] \& C(C(X))
\end{tikzcd}}
sont égales dans\unsure{.} $(\mathcal M_{/S})(W_S^{\mathcal L})^{-1}$. Comme par hyp. on a $W_S^{\mathcal L}=W_S$ dans $\underline{\unsure{\mathrm{Fb}|\mathrm{Fib}}}\,S$ (Hyp 1), et que $\deleted{\unsure{B}}C(X),C(C(X))\in\mathrm{Ob}\,\underline{\mathrm{Fib}}\,S$,
% archive sheet number: 157

\scanpage{159}
\mspage{104}
il suffit de prouver que les deux flèches sont égales \uline{dans} $(\underline{\mathrm{Fib}}\,S)W_S^{-1}$. Cet énoncé est indépendant de la donnée de $\mathcal L$, et on a envie de le prouver.

Peut-être y a-t-il moyen d'\unsure{expliciter|exploiter} un langage géométrique, pour s'y reconnaître. \deleted{Pour \unsure{un objet}} Identifions les objets \inserted{$X$} de $\mathcal M$ aux foncteurs $\mathcal M^\circ\longrightarrow(\mathrm{Ens})$ qu'ils définissent,
\[
X\longmapsto\bigl(T\longmapsto X(T)=\mathrm{Hom}_{\mathcal M}(T,X)\bigr)\text{.}
\]
Alors, pour tt $T$
\[
\mathrm{Ch}(X)(T)
\]
est un ensemble, appelé l'\uline{\deleted{\unsure{ens}}ensemble des chemins dans} $X(T)$. Il s'envoie (par $p_X$) fonctoriellement \inserted{(en $T$)} dans $(X\times X)(T)=X(T)\times X(T)$, i.e. chaque chemin \inserted{$\unsure{c|\gamma}$} dans $X(T)$ \unsure{à|a} une origine \inserted{$\sigma(c)$}, et un but $\beta(c)$. De plus, par $\alpha_X$, $X(T)$ s'\unsure{envoit|envoie} (fonctoriellement en $T$) sur l'ens. des chemins de $X(T)$ : pour chaque élément \inserted{$x$} de $X(T)$, il y a donc un chemin (noté \deleted{$\unsure{\mathrm{id}}$}$\unsure{1}_x$ ou $\alpha_x$) \deleted{\unsure{de}} bien défini de source et but $x$, qu'on appelle le chemin \unsure{trivial|trivide}, ou vide. Ceci dit,
% archive sheet number: 158

\scanpage{160}
\mspage{105}
pour $X$ dans $\deleted{\unsure{C}}\mathcal M_{/S}$, $C(X)$ représente le \inserted{\unsure{contra.}}foncteur suivant dans $\mathcal M\deleted{_{/S}}$ : pour tout $T$ dans $\mathcal M$,
\[
\begin{aligned}
C(X)(T)=\{(x,c)\mid x\in X(T),\quad &\unsure{c}:f_X(x)\xrightarrow{c}\unsure{s}\ \text{un}\\
&\text{chemin dans }S(T),\\
&\text{de source } f_X(x)\}
\end{aligned}
\]
où $f_X:X\longrightarrow S$ est le morphisme structural. Les ``calculs'' \unsure{faits} il y a deux (?) jours montrent qu'on a de $\eqsim$ une interprétation
\[
\begin{aligned}
C(C(X))(T)=\{(x,c,c')\mid x\in X(T),\quad &f_X(x)\mathrel{\overset{c}{\rule{1.6em}{.4pt}}}s\xrightarrow{c'}\unsure{t}\deleted{\unsure{\}}}\\
&c,\text{ un chemin de source }f_X(x),\ c'\\
&\text{un chemin de source égale}\\
&\text{à }\beta(c)\ \}
\end{aligned}
\]
et les applications $i_{C(X)}$ et $\deleted{\unsure{e}}C(i_X)$ s'identifient respectivement à
\[
\begin{aligned}
(x,c)&\longmapsto(x,1_{f(x)},c)\deleted{\unsure{f(x)}}\inserted{\unsure{\text{i.e.}|\mathrm{id}}} && f(x)\xrightarrow{1_x}f(x)\xrightarrow{c}s\\
\deleted{\unsure{\&}}\ (x,c)&\longmapsto(x,c,1_{\beta(c)}) && f(x)\longrightarrow s\xrightarrow{1_s}s\text{.}
\end{aligned}
\]
Pour les comparer, on a une envie irrésistible d'écrire
\[
c\circ1_x=1_s\circ c=c\text{,}
\]
mais en a-t-on le droit ? Il faudrait pouvoir composer les chemins, et c'est très, très difficile !

Mais on peut toujours axiomatiser. On suppose donc que pour $\deleted{\unsure{C}}X$ dans $\mathcal M$, $\mathrm{Ch}(X)$
% archive sheet number: 159

\scanpage{161}
\mspage{106}
\unsure{n'est|n'ait} pas seulement une structure de \inserted{(soi-disant)} ``chemins'' \inserted{dans $X$} amorphe, \deleted{\unsure{avec}} seulement la donnée de la source et du but des chemins et des chemins triviaux (ces structures \inserted{sur $\mathrm{Ch}(X)$} dépendant fonctoriellement en $X$), mais qu'on a une composition des chemins, quand \deleted{\unsure{ceux}} ceux-ci sont ``composables'', i.e. une \deleted{\unsure{application}} flèche d'un produit fibré
\[
(c,c')\longmapsto c'\circ c:(\mathrm{Ch}(X),\beta_X)\times_X(\mathrm{Ch}(X),\sigma_X)\xrightarrow{\pi_X}\mathrm{Ch}(X)\text{,}
\]
i.e. on peut composer \unsure{ses|des} chemins sur \deleted{\unsure{X}} $X(T)$
\diagram{\begin{tikzcd}[ampersand replacement=\&]
x\arrow[r,"c"] \& y\arrow[r,"c'"] \& z
\end{tikzcd}}
ssi $\mathrm{but}\ c=\mathrm{source}\ c'$, et de plus on exige bien sûr que $\mathrm{source}\,(c'\circ c)=\mathrm{source}\ c$, $\mathrm{but}\,(c'\circ c)=\mathrm{but}\ c'$, ce qui s'exprime par un diagramme commutatif
\diagram{\begin{tikzcd}[ampersand replacement=\&]
(\mathrm{Ch}(X),\beta_X)\times_X(\mathrm{Ch}(X),\sigma_X)\arrow[rr]\arrow[dr] \& {} \& \mathrm{Ch}(X)\arrow[dl,"p_X"]\\
{} \& X\times X \& \text{.}
\end{tikzcd}}

La pensée naturelle, c'est qu'avec cette composition et ces ``\unsure{identités|unités}'', les $X(T)$, $\mathrm{Ch}(T)$ deviennent les objets et les flèches d'une \uline{catégorie} --- i.e. que $(X,\mathrm{Ch}(X),\deleted{\unsure{\pi_X}}\deleted{\unsure{\sigma}}\alpha_X,p_X,\pi_X)$ définissent une structure de ``catégorie dans $\mathcal M$''. Mais pour ce qu'on veut faire \inserted{maintenant} on n'a sans doute pas besoin de tout, \deleted{\unsure{de}} \unsure{cela|sla}.
% archive sheet number: 160

\scanpage{162}
\mspage{107}
de l'axiome pertinent d'associativité de la composition
et de celui sur les unités, mais seulement que
les $1_x$ soient des unités bilatères pour la
composition.

En supposant donné \unsure{ce|un} supplément de
structure sur le foncteur $\mathrm{Ch}\deleted{\unsure{(X)|(\ )}}$, \deleted{\illegible{2}} et
revenant à la considération du foncteur
associé $C:\mathcal M_{/S}\longrightarrow\underline{\mathrm{Fib}}\,S$,
on trouve un\deleted{\unsure{e} application} morphisme canonique
\[
C(C(X))\xrightarrow{\gamma_X}C(X)
\]
par
\[
(x,c,c')\longmapsto(x,c'\circ c),
\]
et l'axiome des unités bilatères nous
assure aussitôt que
\[
\unsure{(*)|(a)|(\alpha)}\qquad\qquad \gamma_X\circ\deleted{\unsure{p}}i_{C(X)}=\gamma_X\circ C(i_X)=\mathrm{id}_{C(X)}.
\]
Comme $i_{C(X)}\in W$, cela prouve déjà que
$\gamma_X\in W_{\deleted{S}}$, donc $\gamma_X\in W_S$. Donc $\gamma_X$ devient
un iso dans $(\underline{\mathrm{Fib}}\,S)W_S^{-1}$, et \deleted{\unsure{on aura}}\inserted{la}
relation $\unsure{(*)|(a)|(\alpha)}$ implique bien $i_{C(X)}=C(i_X)$
dans ce localisé, q.e.d.

Je vais à nouveau résumer ce que j'ai
dégagé.
% archive sheet number: 161
